\SourceLeaf{231}%
\bigskip\Needspace{9\baselineskip}
\pdfbookmark[1]{Pallas Olbersiana. October 1802}{pallas180210}
\begin{center}
PALLAS OLBERSIANA.\par
\medskip\hrule\medskip
Monatliche Correspondenz. Vol. VI. pp.\ 390—396. October 1802.\par
\medskip\hrule
\end{center}\medskip
—\quad —\quad —

Dr.\ \textsc{Gauss}, who had improved his elements of the orbit of Pallas for the third time, has now ventured a \textit{fourth} correction. These new elements are derived from our first observations at \textit{Seeberg}, Dr.\ \textsc{Maskelyne}’s last meridian observations at \textit{Greenwich}, and Dr.\ \textsc{Olbers}’s observation of 8 July. Dr.\ \textsc{Gauss} is still unwilling to present them as more reliable than elements III, since he expects a much more accurate correction from the later Milan observations. Here, however, are the elements. [The same as those given above in the letter to \textsc{Bode} of 10 September, p.\ 223, also for the epoch of 31 March 1802 at noon at Seeberg, with “for the epoch, fixed relative to the stars” added for the aphelion and the node.]

Using these elements, Dr.\ \textsc{Gauss} calculated the following ephemeris for the planet’s motion next year. It is useful to publish it in good time. In the meantime, in the region through which \textit{Pallas} will pass in 1803, one can determine the positions of doubtful stars more accurately and get to know and find one’s way around this part of the sky in advance. That will make the planet easier for us to find. The last column of the ephemeris contains the planet’s brightness, with its brightness on 8 July of this year taken as the unit. On the same scale, it was:
\begin{center}
\begin{tabular}{r l l r}
on 4 April&1802&$=$&$4.23$\\
4 May&—&&$2.75$\\
16 July&—&&$0.90$\\
26 July&—&&$0.80$\\
7 Aug.&—&&$0.70$\\
24 Aug.&—&&$0.60$
\end{tabular}
\end{center}
The size of the phase, the altitude above the horizon and the interval from twilight are left out of account here. The first circumstance is entirely negligible. The other two depend on the observing sites and the more powerful telescopes. The greatest brightness \textit{Pallas} will reach next year will be $=0.656$. It had the same brightness this year, 1802, on 14 August. \textsc{Oriani} still \textit{observed} the planet on 8 August and barely \textit{saw} it on the 17th and 18th of that month. Twilight and the planet’s low position probably had much to do with the difficulty. We can therefore still entertain the pleasant hope that we will catch sight of \textit{Pallas} next year and be able to observe its opposition to the Sun. The planet will probably already have been found before then (30 June 1803), for its visibility or brightness already begins to diminish again at that encounter. As for meridian observations, we will probably have little hope of observing this planet before the latter half of May or the beginning of June, at least in our latitudes. [The ephemeris runs from 4 February to 28 June, midnight at Seeberg, 1803. For every third day it gives the R.A. and declination to the nearest minute, the distance from Earth to two decimal places and the brightness to three.]

Even if Dr.\ \textsc{Gauss} succeeds in improving the above elements further with the Italian observation, as is most likely, there will be no need to repeat the laborious
\SourceLeaf{232}%
calculation of an entirely new ephemeris on that basis. It will then be enough to recalculate \textit{some} positions of the planet, since the difference cannot be very considerable and can grow regularly only in slow stages. Dr.\ \textsc{Gauss} has already tried calculating some positions of \textit{Pallas} from his elements III, and the difference is hardly so considerable that it should hinder the recovery of the planet, as the accompanying comparison shows. [The results for elements V have been added here from the manuscript.]
\begin{center}
\begingroup\small\setlength{\tabcolsep}{2.5pt}
\begin{tabular}{l|r|r|r|r|r|r|r|r}
&\multicolumn{2}{c|}{R.A. of $\PallasSymbol$}&&\multicolumn{2}{c|}{Decl. of $\PallasSymbol$}&&R.A.&Decl.\\
1803&III. Elem.&IV. Elem.&Diff.&III. Elem.&IV. Elem.&Diff.&V. Elem.&V. Elem.\\\hline
4 Feb.&$267^0\ 35'$&$267^0\ 48'$&$13'$&$5^0\ 41'$&$5^0\ 38'$&$3'$&$267^0\ 45'$&$5^0\ 39'$\\
28 June&$275\ 45$&$276\ 17$&$32$&$23\ 12$&$23\ 10$&$2$&$276\ 9$&$23\ 11$
\end{tabular}\endgroup
\end{center}
According to elements III, the opposition of this planet falls on the morning of 30 June in civil time, and according to elements IV on the afternoon of the same day. The difference in declination is entirely insignificant. According to V. 30 June. $2^{\mathrm h}\ 53^{\mathrm m}$ In $277^0\ 45'$ long.
\begin{flushright}[\textsc{von Zach.}]\end{flushright}

